\documentclass[12pt, a4paper]{article}

% ============================================================
%  Paquetes
% ============================================================
\usepackage[utf8]{inputenc}
\usepackage[T1]{fontenc}
\usepackage[spanish]{babel}
\usepackage{amsmath, amssymb}
\usepackage[ruled, vlined, linesnumbered, spanish]{algorithm2e}
\usepackage{geometry}
\usepackage{hyperref}

\geometry{margin=2.5cm}

% ============================================================
%  Metadatos del documento
% ============================================================
\title{Cuckoo Search Discreto para CARP --- Pseudocódigo}
\author{}
\date{}

% ============================================================
\begin{document}
\maketitle

% ------------------------------------------------------------
\section{Introducción y analogía biológica}
% ------------------------------------------------------------

El algoritmo \textbf{Cuckoo Search} (Yang \& Deb, 2009) se inspira en el
parasitismo de nidificación del cucú: el pájaro pone sus huevos en los nidos de
otras aves y, si el huevo no es detectado como intruso, el polluelo es criado por
el ave huésped. En términos de optimización:

\begin{itemize}
    \item Cada \textbf{nido} es una solución candidata (un conjunto de rutas CARP).
    \item Cada \textbf{huevo de cucú} es una nueva solución generada mediante un
          \emph{vuelo de Lévy discreto} desde la solución de un nido.
    \item Si el huevo es \emph{mejor} que el nido donde aterriza, lo reemplaza
          (\emph{competencia greedy}).
    \item Con probabilidad $p_a$, los \emph{peores nidos} son abandonados y
          reemplazados por nuevas soluciones generadas en torno al mejor nido
          (\emph{diversificación}).
\end{itemize}

La diferencia con ABC es que Cuckoo Search genera sus soluciones candidatas
mediante el vuelo de Lévy, que combina muchos pasos cortos con saltos
ocasionalmente largos, en lugar de la perturbación unitaria de las empleadas
o la selección por ruleta de las observadoras.

% ------------------------------------------------------------
\section{Notación}
% ------------------------------------------------------------

\begin{itemize}
    \item $G = (V, E, R)$: grafo de la instancia CARP. $R \subseteq E$ son los
          arcos requeridos.
    \item $N = \{s_1, \ldots, s_{K}\}$: conjunto de nidos (soluciones activas),
          con $K = \mathit{num\_nidos}$.
    \item $s^{*}$: mejor solución encontrada durante toda la corrida.
    \item $c(s)$: costo puro de la solución $s$.
    \item $\mathrm{viol}(s)$: exceso total de demanda respecto a la capacidad.
    \item $f(s) = c(s) + \lambda \cdot \mathrm{viol}(s)$: objetivo penalizado
          ($\lambda > 0$).
    \item $p_a \in (0,1)$: fracción de peores nidos abandonados por iteración.
    \item $\beta > 0$: parámetro de forma del vuelo de Lévy ($\beta = 1.5$
          clásico).
    \item $L$: número de pasos del vuelo de Lévy discreto (variable aleatoria
          acotada a $[1, 12]$, ver sección~\ref{sec:levy}).
    \item $P_{\mathrm{base}}$: escala base de los pasos de Lévy.
    \item $\mathcal{O}_{\mathrm{intra}}$, $\mathcal{O}_{\mathrm{inter}}$:
          conjuntos de operadores intra-ruta e inter-ruta (9 operadores en total).
    \item $p_{\mathrm{inter}} \in [0,1]$: P(elegir inter-ruta) cuando la solución
          es factible.
    \item $p_{\mathrm{efec}}$: probabilidad efectiva de elegir inter-ruta tras el
          piso dinámico bajo violación.
    \item $n = |R|$: número de tareas requeridas (variable de escala
          \emph{instance-aware}).
\end{itemize}

\medskip
\noindent
\textbf{Calibración \emph{instance-aware} (valores por defecto si no se
proporcionan):}
\[
\mathit{iteraciones} = \max(200,\; 20 \cdot n), \quad
K = \max(10,\; \lceil 2\sqrt{n} \rceil), \quad
P_{\mathrm{base}} = \max\!\left(3,\; \left\lfloor\tfrac{\sqrt{n}}{2}\right\rfloor\right), \quad
M = \max(50,\; 3 \cdot n).
\]

\medskip
\noindent
\textbf{Sesgo inter/intra dinámico (compartido con SA, TS simple, RTS y ABC simple).}
Se calcula la violación media $\bar{v} = \tfrac{1}{K}\sum_{i=1}^{K}\mathrm{viol}(s_i)$
y se aplica:
\[
p_{\mathrm{efec}} \;=\;
\begin{cases}
  \max(p_{\mathrm{inter}},\; 0.8) & \text{si } \bar{v} > 0, \\
  p_{\mathrm{inter}}              & \text{si } \bar{v} = 0.
\end{cases}
\]

% ------------------------------------------------------------
\section{Vuelo de Lévy discreto}
\label{sec:levy}
% ------------------------------------------------------------

En el espacio continuo original, el vuelo de Lévy produce pasos de longitud con
distribución de cola pesada. En el espacio discreto de rutas CARP se aproxima
por el \emph{número de perturbaciones consecutivas} $L$:

\[
x \sim |\mathcal{N}(0,\,1)|, \qquad
L = \min\!\left(12,\; 1 + \left\lfloor x^{1/\beta} \cdot P_{\mathrm{base}} \right\rfloor\right).
\]

Un $x$ grande (cola de la distribución normal) produce un $L$ grande y por tanto
un salto largo en el espacio de soluciones. La cota $L \leq 12$ evita vuelos
excesivamente largos que degradarían la calidad de la exploración.

Cada uno de los $L$ pasos es una perturbación local: aplicar un operador
$\mathcal{O}_{\mathrm{grupo}}$ (elegido por el sesgo inter/intra) a la solución
actual y obtener la solución resultante. El operador concreto dentro del grupo
se selecciona de forma uniforme.

% ------------------------------------------------------------
\section{Pseudocódigos}
% ------------------------------------------------------------

% -------- Lévy discreto (auxiliar) --------------------------
\begin{algorithm}[H]
\DontPrintSemicolon
\caption{VueloLévyDiscreto($s_{\mathrm{base}},\, P_{\mathrm{base}},\, \beta,\,
         \mathcal{O}_{\mathrm{grupo}}$)}
\label{alg:levy}

\KwIn{Solución base $s_{\mathrm{base}}$;\; escala $P_{\mathrm{base}}$;\;
      parámetro $\beta$;\; grupo de operadores $\mathcal{O}_{\mathrm{grupo}}$}
\KwOut{Solución cuckoo $s'$;\; lista de movimientos aplicados $\mathbf{m}$}

\BlankLine
$x \leftarrow |{\cal N}(0,1)|$
    \tcp{valor absoluto de una muestra normal estándar}
$L \leftarrow \min\!\left(12,\; 1 + \lfloor x^{1/\beta} \cdot
    \max(1, P_{\mathrm{base}}) \rfloor\right)$
    \tcp{número de pasos del vuelo; acotado en $[1,12]$}
$s' \leftarrow$ copia de $s_{\mathrm{base}}$ \;
$\mathbf{m} \leftarrow [\,]$ \;
\For{$\ell \leftarrow 1$ \KwTo $L$}{
    $(s',\, m_\ell) \leftarrow \textsc{GenerarVecino}(s',\, \mathcal{O}_{\mathrm{grupo}})$
        \tcp{perturbación local (operador uniforme dentro del grupo)}
    $\mathbf{m} \leftarrow \mathbf{m} + [m_\ell]$ \;
}
\Return $(s',\, \mathbf{m})$
\end{algorithm}

\bigskip

% -------- Cuckoo Search completo ----------------------------
\begin{algorithm}[H]
\DontPrintSemicolon
\caption{Cuckoo Search Discreto para CARP (Yang \& Deb 2009 adaptado)}
\label{alg:cuckoo}

\KwIn{Instancia CARP $G$;\; solución inicial $s_0$;\;
      $p_a$;\; $\beta$;\; $p_{\mathrm{inter}}$;\; $\lambda$;\;
      opcionalmente: $K$, $P_{\mathrm{base}}$, $\mathit{iteraciones}$,
      $M$ (todos instance-aware si se omiten)}
\KwOut{Mejor solución encontrada $s^{*}$}

\BlankLine
\tcp{Calibración instance-aware}
$n \leftarrow |R|$ \;
$\mathit{iteraciones} \leftarrow \max(200,\; 20 \cdot n)$ \;
$K            \leftarrow \max(10,\; \lceil 2\sqrt{n} \rceil)$ \;
$P_{\mathrm{base}} \leftarrow \max(3,\; \lfloor\sqrt{n}/2\rfloor)$ \;
$M            \leftarrow \max(50,\; 3 \cdot n)$
    \tcp{máx. iteraciones sin mejora}

\BlankLine
\tcp{Precómputo UNA VEZ de las particiones intra/inter}
Construir $\mathcal{O}_{\mathrm{intra}}$, $\mathcal{O}_{\mathrm{inter}}$,
$\mathcal{O}_{\mathrm{fallback}}$ a partir de los operadores activos \;

\BlankLine
\tcp{Inicialización de nidos}
$N[0] \leftarrow s_0$ \;
\For{$i \leftarrow 1$ \KwTo $K - 1$}{
    $N[i] \leftarrow$ vecino aleatorio de $N[0]$ \;
}
$s^{*} \leftarrow \arg\min_{i} f(N[i])$ \;
$\mathit{iter\_sin\_mejora} \leftarrow 0$ \;

\BlankLine
\For{$\mathit{ciclo} \leftarrow 1$ \KwTo $\mathit{iteraciones}$}{

    \tcp{Criterio de parada anticipada por estancamiento}
    \If{$\mathit{iter\_sin\_mejora} \geq M$}{
        \textbf{break} \;
    }
    $f^{*}_{\mathrm{inicio}} \leftarrow f(s^{*})$ \tcp{snapshot para detectar mejora}

    \BlankLine
    \tcp{Violación media y $p_{\mathrm{efec}}$ para este ciclo}
    $\bar{v} \leftarrow \tfrac{1}{K}\sum_{i=1}^{K} \mathrm{viol}(N[i])$ \;
    \lIf{$\bar{v} > 0$}{$p_{\mathrm{efec}} \leftarrow \max(p_{\mathrm{inter}},\; 0.8)$}
    \lElse{$p_{\mathrm{efec}} \leftarrow p_{\mathrm{inter}}$}

    \BlankLine
    \tcp{\textbf{Paso 1 — Generación de cuckoos mediante vuelo de Lévy}}
    \For{$i \leftarrow 0$ \KwTo $K - 1$}{
        $\mathcal{O}_{\mathrm{grupo}} \leftarrow$ \textsc{ElegirGrupo}($p_{\mathrm{efec}},\,
            \mathrm{viol}(N[i]),\,
            \mathcal{O}_{\mathrm{inter}},\, \mathcal{O}_{\mathrm{intra}}$) \;
        $(\mathit{ck}[i],\, \mathbf{m}[i]) \leftarrow$ \textsc{VueloLévyDiscreto}($N[i],\,
            P_{\mathrm{base}},\, \beta,\, \mathcal{O}_{\mathrm{grupo}}$) \;
        \ForEach{$m \in \mathbf{m}[i]$}{\textsc{Proponer}($m$)} \;
    }

    \BlankLine
    \tcp{Evaluación vectorizada del lote de cuckoos (GPU si disponible)}
    $\{f(\mathit{ck}[i])\}_{i=0}^{K-1} \leftarrow$ \textsc{EvaluarLote}($\{\mathit{ck}[i]\}$) \;

    \BlankLine
    \tcp{\textbf{Paso 2 — Competencia: cuckoo vs nido aleatorio}}
    \For{$i \leftarrow 0$ \KwTo $K - 1$}{
        $j \leftarrow \textsc{UniformeEntero}(0,\, K-1)$
            \tcp{nido rival elegido al azar}
        \If{$f(\mathit{ck}[i]) < f(N[j]) - \varepsilon$}{
            $N[j] \leftarrow \mathit{ck}[i]$
                \tcp{el cuckoo reemplaza al nido rival}
            \textsc{Aceptar}($\mathbf{m}[i][-1]$) \;
            \tcp{Detección inmediata de mejora global (corrección del bug de imputación)}
            \If{$f(\mathit{ck}[i]) < f(s^{*}) - \varepsilon$}{
                $s^{*} \leftarrow \mathit{ck}[i]$ \;
                \textsc{RegistrarMejora}($\mathbf{m}[i][-1]$)
                    \tcp{operador del último paso del vuelo Lévy de ESTE cuckoo}
            }
        }
    }

    \BlankLine
    \tcp{\textbf{Paso 3 — Abandono de los peores nidos}}
    $n_a \leftarrow \max(1,\; \lfloor p_a \cdot K \rfloor)$
        \tcp{cantidad de nidos a abandonar}
    $\mathit{peores} \leftarrow$ índices de los $n_a$ nidos con mayor $f(\cdot)$ \;
    $j^{*} \leftarrow \arg\min_{i} f(N[i])$
        \tcp{índice del mejor nido}
    $\mathcal{O}_{\mathrm{ab}} \leftarrow$ \textsc{ElegirGrupo}($p_{\mathrm{efec}},\,
        \mathrm{viol}(N[j^{*}]),\,
        \mathcal{O}_{\mathrm{inter}},\, \mathcal{O}_{\mathrm{intra}}$) \;
    \ForEach{$k \in \mathit{peores}$}{
        $(N[k],\, \mathbf{m}_{\mathrm{ab}}) \leftarrow$ \textsc{VueloLévyDiscreto}($N[j^{*}],\,
            P_{\mathrm{base}},\, \beta,\, \mathcal{O}_{\mathrm{ab}}$)
            \tcp{vuelo Lévy desde el mejor nido actual}
        \ForEach{$m \in \mathbf{m}_{\mathrm{ab}}$}{\textsc{Proponer}($m$)} \;
        \tcp{(evaluación incluida en el bloque final)}
    }

    \BlankLine
    \tcp{Evaluación vectorizada de los nidos recién reemplazados}
    Evaluar $f(N[k])$ para todo $k \in \mathit{peores}$ \;
    \ForEach{$k \in \mathit{peores}$}{
        \If{$f(N[k]) < f(s^{*}) - \varepsilon$}{
            $s^{*} \leftarrow N[k]$ \tcp{el abandono puede también mejorar el global}
        }
    }

    \BlankLine
    \tcp{Actualización del contador de estancamiento}
    \lIf{$f(s^{*}) < f^{*}_{\mathrm{inicio}} - \varepsilon$}{%
        $\mathit{iter\_sin\_mejora} \leftarrow 0$}
    \lElse{%
        $\mathit{iter\_sin\_mejora} \leftarrow \mathit{iter\_sin\_mejora} + 1$}
}

\BlankLine
\Return $s^{*}$
\end{algorithm}

% ------------------------------------------------------------
\section{Notas de implementación}
% ------------------------------------------------------------

\paragraph{Evaluación vectorizada (GPU).}
En el Paso 1, los $K$ cuckoos se evalúan en lote con
\texttt{costo\_lote\_penalizado\_ids} (NumPy en CPU; CuPy en GPU si
\texttt{usar\_gpu=True}). Los nidos reemplazados en el Paso 3 también se evalúan
en lote. Esto concentra el trabajo vectorizable en dos llamadas por iteración,
minimizando el overhead de Python puro.

\paragraph{Corrección del bug de imputación.}
En la versión anterior, la detección de mejora global se hacía comparando el
mejor antes y después de la iteración completa, imputando el crédito al
\emph{último movimiento aceptado} de la iteración (que podía no ser el operador
responsable). En esta versión, la llamada a \textsc{RegistrarMejora}$(m)$ ocurre
dentro del mismo \texttt{if} que detecta $f(\mathit{ck}[i]) < f(s^{*})$, con el
operador del último paso del vuelo de Lévy de \emph{ese} cuckoo específico.
Consecuencia: $\sum_{\mathrm{op}}\mathit{mejoraron}[\mathrm{op}] = \mathit{mejoras}$.

\paragraph{Inicialización de nidos.}
El nido 0 se siembra con la mejor solución inicial disponible. Los nidos
$1 \ldots K-1$ se generan como vecinos aleatorios del nido 0 (igual estrategia
que en ABC simple), lo que ofrece diversidad inmediata sin desperdiciar la
información de la solución de partida.

\paragraph{Diferencia con ABC simple en los scouts.}
En Cuckoo Search, el paso de abandono (Paso 3) reemplaza los peores nidos con
un vuelo de Lévy \emph{desde el mejor nido actual} (exploración dirigida). En
ABC simple, los scouts reemplazan fuentes agotadas con soluciones completamente
aleatorias (exploración pura). Cuckoo Search tiende a explorar en torno a la
región de mejor calidad conocida; ABC simple mantiene mayor diversificación global.

\paragraph{Parámetro $\beta$.}
El valor clásico de Yang \& Deb es $\beta = 1.5$. Valores menores (ej.\ $\beta =
1.2$) generan distribuciones con cola más pesada: más saltos largos, mayor
exploración pero menor precisión local. Valores mayores (ej.\ $\beta = 2.0$)
producen pasos más cortos, favoreciendo la explotación. En MetaCARP se usa el
valor clásico como default.

\end{document}
